\documentclass[12pt,a4paper]{article}
\usepackage[utf8]{inputenc}
\usepackage[english]{babel}
\usepackage{geometry}
\usepackage{amsmath}
\usepackage{graphicx}
\usepackage{hyperref}
\usepackage{booktabs}

\geometry{margin=1in}

\title{Classification of BBC Articles into Categories}
\author{Krzysztof Czechowicz \and Jakub Błażowski \and Maciej Pawłowski}
\date{November 24, 2025}

\begin{document}

\maketitle

\begin{abstract}
Different supervised machine learning models were compared for classifying BBC articles. All articles were scraped from the official BBC website. Several approaches were tested, from traditional algorithms like Naive Bayes to deep learning models. The goal was to see which methods work best for categorizing news and sports content.
\end{abstract}

\section{Introduction}

The work was split into two parts. First, sports articles were classified into specific sports like football, rugby, tennis, boxing, and a few others. Then general news classification was done, putting articles into categories such as politics, sport, technology, business, health, culture, and world news. For both tasks, different models were tried: Naive Bayes, Support Vector Machines, and Convolutional Neural Networks. The goal was to see which ones give the best results.

\section{Classification of Sports Articles}

\subsection{Dataset}

The training data was split in an 80/20 ratio. The number of articles and their categories are shown in Table~\ref{tab:sports_dataset}.

\begin{table}[h]
\centering
\begin{tabular}{lcc}
\toprule
Category & Training Set & Test Set \\
\midrule
Rugby Union & 320 & 80 \\
Football & 320 & 80 \\
Formula 1 & 320 & 80 \\
Cricket & 320 & 80 \\
Boxing & 307 & 77 \\
Rugby League & 267 & 67 \\
Tennis & 213 & 53 \\
American Football & 199 & 50 \\
Mixed Martial Arts & 167 & 42 \\
\bottomrule
\end{tabular}
\caption{Distribution of sports articles in training and test sets}
\label{tab:sports_dataset}
\end{table}

\subsection{Naive Bayes}

CountVectorizer was used to convert text into numerical features. It creates a matrix where each row represents a document and each column represents a word, counting how many times each word appears. The parameter \texttt{min\_df = 2} means that words appearing in fewer than 2 documents are ignored. \texttt{max\_features = 10000} limits the vocabulary to the 10,000 most frequent words. Smoothing parameter was set to 1.0.

The model achieved an overall accuracy of 99.34\%. Detailed results are shown in Table~\ref{tab:sports_nb}.

\begin{table}[h]
\centering
\begin{tabular}{lccc}
\toprule
Category & Precision & Recall & F1-Score \\
\midrule
American Football & 1.0000 & 1.0000 & 1.0000 \\
Boxing & 1.0000 & 1.0000 & 1.0000 \\
Cricket & 1.0000 & 1.0000 & 1.0000 \\
Football & 1.0000 & 0.9625 & 0.9809 \\
Formula 1 & 1.0000 & 1.0000 & 1.0000 \\
Mixed Martial Arts & 1.0000 & 1.0000 & 1.0000 \\
Rugby League & 1.0000 & 0.9851 & 0.9925 \\
Rugby Union & 0.9524 & 1.0000 & 0.9756 \\
Tennis & 1.0000 & 1.0000 & 1.0000 \\
\bottomrule
\end{tabular}
\caption{Naive Bayes classification results for sports articles}
\label{tab:sports_nb}
\end{table}

\subsection{LinearSVC}

Two different vectorization approaches were tested with LinearSVC. Both used \texttt{min\_df = 2} and \texttt{max\_features = 10000}. The LinearSVC classifier was configured with \texttt{random\_state = 42} and \texttt{max\_iter = 2000}.

First, CountVectorizer was used, achieving an overall accuracy of 98.69\%. Results are shown in Table~\ref{tab:sports_svc_count}.

\begin{table}[h]
\centering
\begin{tabular}{lccc}
\toprule
Category & Precision & Recall & F1-Score \\
\midrule
American Football & 1.0000 & 1.0000 & 1.0000 \\
Boxing & 0.9872 & 1.0000 & 0.9935 \\
Cricket & 0.9877 & 1.0000 & 0.9938 \\
Football & 0.9630 & 0.9750 & 0.9689 \\
Formula 1 & 1.0000 & 1.0000 & 1.0000 \\
Mixed Martial Arts & 1.0000 & 0.9524 & 0.9756 \\
Rugby League & 1.0000 & 0.9701 & 0.9848 \\
Rugby Union & 0.9634 & 0.9875 & 0.9753 \\
Tennis & 1.0000 & 0.9811 & 0.9905 \\
\bottomrule
\end{tabular}
\caption{LinearSVC with CountVectorizer classification results for sports articles}
\label{tab:sports_svc_count}
\end{table}

Then, TF-IDF Vectorizer was tested. TF-IDF (Term Frequency-Inverse Document Frequency) weights words by how frequently they appear in a document relative to how rare they are across all documents. This approach achieved an overall accuracy of 99.84\%, which is higher than the Naive Bayes baseline. Detailed results are shown in Table~\ref{tab:sports_svc_tfidf}.

\begin{table}[h]
\centering
\begin{tabular}{lccc}
\toprule
Category & Precision & Recall & F1-Score \\
\midrule
American Football & 1.0000 & 1.0000 & 1.0000 \\
Boxing & 1.0000 & 1.0000 & 1.0000 \\
Cricket & 1.0000 & 1.0000 & 1.0000 \\
Football & 1.0000 & 1.0000 & 1.0000 \\
Formula 1 & 1.0000 & 1.0000 & 1.0000 \\
Mixed Martial Arts & 1.0000 & 1.0000 & 1.0000 \\
Rugby League & 1.0000 & 0.9851 & 0.9925 \\
Rugby Union & 0.9877 & 1.0000 & 0.9938 \\
Tennis & 1.0000 & 1.0000 & 1.0000 \\
\bottomrule
\end{tabular}
\caption{LinearSVC with TF-IDF Vectorizer classification results for sports articles}
\label{tab:sports_svc_tfidf}
\end{table}

\subsection{CNN (Conv1D)}

A Convolutional Neural Network with 1D convolutions was used. The Keras Tokenizer was applied with \texttt{max\_words = 10000} and \texttt{max\_len = 500} for sequence padding. The model architecture included an Embedding layer with dimension 128, a Conv1D layer with 128 filters and ReLU activation, followed by GlobalMaxPooling1D, a Dense layer with 64 units and ReLU activation, and a final Dense layer with softmax activation for classification.

The model achieved an overall accuracy of 97.21\%. Detailed results are shown in Table~\ref{tab:sports_cnn}.

\begin{table}[h]
\centering
\begin{tabular}{lccc}
\toprule
Category & Precision & Recall & F1-Score \\
\midrule
American Football & 0.9608 & 0.9800 & 0.9703 \\
Boxing & 0.9615 & 0.9740 & 0.9677 \\
Cricket & 0.9753 & 0.9875 & 0.9814 \\
Football & 0.9753 & 0.9875 & 0.9814 \\
Formula 1 & 1.0000 & 0.9875 & 0.9937 \\
Mixed Martial Arts & 0.9318 & 0.9762 & 0.9535 \\
Rugby League & 0.9701 & 0.9701 & 0.9701 \\
Rugby Union & 0.9737 & 0.9250 & 0.9487 \\
Tennis & 0.9808 & 0.9623 & 0.9714 \\
\bottomrule
\end{tabular}
\caption{CNN (Conv1D) classification results for sports articles}
\label{tab:sports_cnn}
\end{table}

\subsection{Key Words}

The most important words that influenced classification decisions were analyzed for both Naive Bayes and LinearSVC models. Table~\ref{tab:sports_keywords} shows the top 3 words for each category.

\begin{table}[h]
\centering
\begin{tabular}{lll}
\toprule
Category & Naive Bayes & LinearSVC \\
\midrule
American Football & quarterback, nfl, chiefs & nfl, bowl, field \\
Boxing & usyk, dubois, fury & boxing, fight, world \\
Cricket & wickets, wicket, cricket & cricket, county, wickets \\
Football & midfielder, rangers, arsenal & use, football, manager \\
Formula 1 & verstappen, norris, mclaren & prix, formula, race \\
Mixed Martial Arts & ufc, pfl, aspinall & ufc, mma, pfl \\
Rugby League & kr, wigan, helens & hull, nrl, super \\
Rugby Union & wru, ospreys, scarlets & rugby, premiership, irish \\
Tennis & sinner, alcaraz, raducanu & open, tennis, wimbledon \\
\bottomrule
\end{tabular}
\caption{Top 3 key words for each category: Naive Bayes vs LinearSVC}
\label{tab:sports_keywords}
\end{table}
\newpage
\subsection{Model Comparison}

All tested models achieved high accuracy scores, as shown in Table~\ref{tab:sports_comparison}. These excellent results can be explained by the nature of the task. Sports articles contain many category-specific keywords that appear frequently and repeatedly. For example, words like "boxing", "fight", and boxer names appear consistently in boxing articles, making classification relatively straightforward for the models.

\begin{table}[h]
\centering
\begin{tabular}{lc}
\toprule
Model & Accuracy \\
\midrule
LinearSVC (TF-IDF) & 99.84\% \\
Naive Bayes & 99.34\% \\
LinearSVC (CountVectorizer) & 98.69\% \\
CNN (Conv1D) & 97.21\% \\
\bottomrule
\end{tabular}
\caption{Accuracy comparison of all models for sports article classification}
\label{tab:sports_comparison}
\end{table}

The best performing model was LinearSVC with TF-IDF vectorization, which considers the weight of individual words and their importance across documents. Naive Bayes came in second place, followed by the standard LinearSVC with CountVectorizer. The CNN model performed worst despite being the most complex and advanced architecture. This can be explained by the fact that simpler models work well when the task relies heavily on keyword presence, while deep learning models may overcomplicate the problem or require more data to learn effectively. For this specific task, the straightforward word-based features were sufficient, and the additional complexity of the neural network did not provide benefits.
\newpage
\section{Classification of News Articles}

\subsection{Dataset}

The training data was split in an 80/20 ratio. The number of articles and their categories are shown in Table~\ref{tab:news_dataset}.

\begin{table}[h]
\centering
\begin{tabular}{lcc}
\toprule
Category & Training Set & Test Set \\
\midrule
World & 800 & 200 \\
Health & 800 & 200 \\
Politics & 800 & 200 \\
Sport & 800 & 200 \\
Business & 800 & 200 \\
Culture & 800 & 200 \\
Technology & 800 & 200 \\
\bottomrule
\end{tabular}
\caption{Distribution of news articles in training and test sets}
\label{tab:news_dataset}
\end{table}

\subsection{Naive Bayes}

Three different vectorization approaches were tested with Naive Bayes classifier. All configurations used the same parameters: \texttt{min\_df = 2}, \texttt{max\_features = 10000}, and smoothing parameter \texttt{alpha = 1.0}.

\subsubsection{Count Vectorizer}

CountVectorizer was used to convert text into numerical features by counting word occurrences. The model achieved an overall accuracy of 86.07\%. Detailed results are shown in Table~\ref{tab:news_nb_count}.

\begin{table}[h]
\centering
\begin{tabular}{lccc}
\toprule
Category & Precision & Recall & F1-Score \\
\midrule
Business & 0.8247 & 0.8000 & 0.8122 \\
Culture & 0.7783 & 0.8950 & 0.8326 \\
Health & 0.8312 & 0.9600 & 0.8910 \\
Politics & 0.9430 & 0.7450 & 0.8324 \\
Technology & 0.7746 & 0.8250 & 0.7990 \\
World & 0.9368 & 0.8150 & 0.8717 \\
Sport & 0.9850 & 0.9850 & 0.9850 \\
\bottomrule
\end{tabular}
\caption{Naive Bayes with CountVectorizer classification results for news articles}
\label{tab:news_nb_count}
\end{table}

\subsubsection{TF-IDF}

TF-IDF Vectorizer was tested with the same parameters. This approach weights words by their frequency in documents relative to their rarity across the corpus. The model achieved an overall accuracy of 85.86\%. Detailed results are shown in Table~\ref{tab:news_nb_tfidf}.

\begin{table}[h]
\centering
\begin{tabular}{lccc}
\toprule
Category & Precision & Recall & F1-Score \\
\midrule
Business & 0.8200 & 0.8200 & 0.8200 \\
Culture & 0.7877 & 0.8350 & 0.8107 \\
Health & 0.8101 & 0.9600 & 0.8787 \\
Politics & 0.9226 & 0.7750 & 0.8424 \\
Technology & 0.8112 & 0.7950 & 0.8030 \\
World & 0.9076 & 0.8350 & 0.8698 \\
Sport & 0.9754 & 0.9900 & 0.9826 \\
\bottomrule
\end{tabular}
\caption{Naive Bayes with TF-IDF Vectorizer classification results for news articles}
\label{tab:news_nb_tfidf}
\end{table}

\subsubsection{Word2Vec}

Word2Vec embeddings were used with Gaussian Naive Bayes classifier. Word2Vec creates dense vector representations of words by learning from context. The model achieved an overall accuracy of 73.43\%. Detailed results are shown in Table~\ref{tab:news_nb_w2v}.

\begin{table}[h]
\centering
\begin{tabular}{lccc}
\toprule
Category & Precision & Recall & F1-Score \\
\midrule
Business & 0.8200 & 0.8200 & 0.8200 \\
Culture & 0.7877 & 0.8350 & 0.8107 \\
Health & 0.8101 & 0.9600 & 0.8787 \\
Politics & 0.9226 & 0.7750 & 0.8424 \\
Technology & 0.8112 & 0.7950 & 0.8030 \\
World & 0.9076 & 0.8350 & 0.8698 \\
Sport & 0.9754 & 0.9900 & 0.9826 \\
\bottomrule
\end{tabular}
\caption{Naive Bayes with Word2Vec embeddings classification results for news articles}
\label{tab:news_nb_w2v}
\end{table}

Table~\ref{tab:news_nb_comparison} shows a comparison of accuracy across different Naive Bayes vectorization approaches.

\begin{table}[h]
\centering
\begin{tabular}{lc}
\toprule
Vectorization Method & Accuracy \\
\midrule
CountVectorizer & 86.07\% \\
TF-IDF & 85.86\% \\
Word2Vec & 73.43\% \\
\bottomrule
\end{tabular}
\caption{Accuracy comparison of Naive Bayes with different vectorization methods}
\label{tab:news_nb_comparison}
\end{table}

CountVectorizer achieved slightly better results than TF-IDF, likely because raw word frequencies provide sufficient discriminative power for this task where categories have distinct vocabularies. Word2Vec performed significantly worse because averaging word embeddings loses important information and Gaussian Naive Bayes may not be well-suited for the distribution of Word2Vec vectors.

\subsection{LinearSVC}

Three different vectorization approaches were tested with LinearSVC classifier. The classifier was configured with \texttt{random\_state = 42} and \texttt{max\_iter = 2000}.

\subsubsection{Count Vectorizer}

CountVectorizer was used with \texttt{min\_df = 2} and \texttt{max\_features = 10000}. The model achieved an overall accuracy of 80.00\%. Detailed results are shown in Table~\ref{tab:news_svc_count}.

\begin{table}[h]
\centering
\begin{tabular}{lccc}
\toprule
Category & Precision & Recall & F1-Score \\
\midrule
Business & 0.7228 & 0.7300 & 0.7264 \\
Culture & 0.7552 & 0.7250 & 0.7398 \\
Health & 0.8365 & 0.8700 & 0.8529 \\
Politics & 0.8034 & 0.7150 & 0.7566 \\
Technology & 0.6849 & 0.7500 & 0.7160 \\
World & 0.8513 & 0.8300 & 0.8405 \\
Sport & 0.9515 & 0.9800 & 0.9655 \\
\bottomrule
\end{tabular}
\caption{LinearSVC with CountVectorizer classification results for news articles}
\label{tab:news_svc_count}
\end{table}

\subsubsection{TF-IDF Vectorizer}

TF-IDF Vectorizer was tested with the same parameters (\texttt{min\_df = 2} and \texttt{max\_features = 10000}). The model achieved an overall accuracy of 86.36\%. Detailed results are shown in Table~\ref{tab:news_svc_tfidf}.

\begin{table}[h]
\centering
\begin{tabular}{lccc}
\toprule
Category & Precision & Recall & F1-Score \\
\midrule
Business & 0.8281 & 0.7950 & 0.8112 \\
Culture & 0.8450 & 0.8450 & 0.8450 \\
Health & 0.8551 & 0.9150 & 0.8841 \\
Politics & 0.8833 & 0.7950 & 0.8368 \\
Technology & 0.7840 & 0.8350 & 0.8087 \\
World & 0.8918 & 0.8650 & 0.8782 \\
Sport & 0.9614 & 0.9950 & 0.9779 \\
\bottomrule
\end{tabular}
\caption{LinearSVC with TF-IDF Vectorizer classification results for news articles}
\label{tab:news_svc_tfidf}
\end{table}

\subsubsection{Word2Vec}

Word2Vec embeddings were used with LinearSVC classifier. The model achieved an overall accuracy of 83.86\%. Detailed results are shown in Table~\ref{tab:news_svc_w2v}.

\begin{table}[h]
\centering
\begin{tabular}{lccc}
\toprule
Category & Precision & Recall & F1-Score \\
\midrule
Business & 0.7861 & 0.6800 & 0.7292 \\
Culture & 0.8308 & 0.8100 & 0.8203 \\
Health & 0.8341 & 0.8800 & 0.8564 \\
Politics & 0.8392 & 0.8350 & 0.8371 \\
Technology & 0.7356 & 0.7650 & 0.7500 \\
World & 0.8708 & 0.9100 & 0.8900 \\
Sport & 0.9659 & 0.9900 & 0.9778 \\
\bottomrule
\end{tabular}
\caption{LinearSVC with Word2Vec embeddings classification results for news articles}
\label{tab:news_svc_w2v}
\end{table}

Table~\ref{tab:news_svc_comparison} shows a comparison of accuracy across different LinearSVC vectorization approaches.

\begin{table}[h]
\centering
\begin{tabular}{lc}
\toprule
Vectorization Method & Accuracy \\
\midrule
TF-IDF & 86.36\% \\
Word2Vec & 83.86\% \\
CountVectorizer & 80.00\% \\
\bottomrule
\end{tabular}
\caption{Accuracy comparison of LinearSVC with different vectorization methods}
\label{tab:news_svc_comparison}
\end{table}

TF-IDF achieved the best results for LinearSVC, as it better captures word importance and relevance across documents, which aligns well with the linear decision boundaries of SVM classifiers.

\newpage
\subsection{CNN (Conv1D)}

A Convolutional Neural Network with 1D convolutions was used for news article classification. The Keras Tokenizer was applied with \texttt{max\_words = 10000} and \texttt{max\_len = 500} for sequence padding. The model architecture included an Embedding layer with dimension 128, a Conv1D layer with 128 filters, kernel size of 5, and ReLU activation, followed by GlobalMaxPooling1D, a Dense layer with 64 units and ReLU activation, Dropout with rate 0.5, and a final Dense layer with softmax activation for classification. The model was compiled with Adam optimizer and categorical crossentropy loss. Early stopping was applied with patience of 3 epochs monitoring validation loss.

Three different embedding approaches were tested: random embeddings, Word2Vec embeddings, and BERT embeddings, each with various dimensions.

\subsubsection{Random Embedding Type}

Random embeddings were tested with different dimensions. The embeddings are randomly initialized and learned during training. Results for different dimensions are shown in Table~\ref{tab:news_cnn_random}.

\begin{table}[h]
\centering
\begin{tabular}{lc}
\toprule
Embedding Dimension & Accuracy \\
\midrule
128 & 80.50\% \\
256 & 79.64\% \\
64 & 79.07\% \\
\bottomrule
\end{tabular}
\caption{CNN with Random Embeddings: Accuracy by dimension}
\label{tab:news_cnn_random}
\end{table}

\subsubsection{Word2Vec Embedding Type}

Word2Vec embeddings were tested with different dimensions. Word2Vec models were trained separately for each dimension, and the resulting embeddings were used to initialize the CNN embedding layer. Results are shown in Table~\ref{tab:news_cnn_w2v}.

\begin{table}[h]
\centering
\begin{tabular}{lc}
\toprule
Embedding Dimension & Accuracy \\
\midrule
50 & 78.00\% \\
200 & 77.93\% \\
100 & 76.21\% \\
\bottomrule
\end{tabular}
\caption{CNN with Word2Vec Embeddings: Accuracy by dimension}
\label{tab:news_cnn_w2v}
\end{table}

\subsubsection{BERT Embedding Type}

BERT embeddings were extracted using the pre-trained \texttt{bert-base-uncased} model. BERT provides token-level embeddings with 768 dimensions, which were then reduced using PCA to test different dimensions. Results are shown in Table~\ref{tab:news_cnn_bert}.

\begin{table}[h]
\centering
\begin{tabular}{lc}
\toprule
Embedding Dimension & Accuracy \\
\midrule
256 & 83.79\% \\
128 & 83.36\% \\
512 & 82.86\% \\
\bottomrule
\end{tabular}
\caption{CNN with BERT Embeddings: Accuracy by dimension}
\label{tab:news_cnn_bert}
\end{table}

\subsection{Summary}

Table~\ref{tab:news_cnn_comparison} shows a comparison of the best accuracy achieved by each embedding type in CNN models.

\begin{table}[h]
\centering
\begin{tabular}{lc}
\toprule
Embedding Type & Best Accuracy \\
\midrule
BERT (256-dim) & 83.79\% \\
Random (128-dim) & 80.50\% \\
Word2Vec (50-dim) & 78.00\% \\
\bottomrule
\end{tabular}
\caption{Best accuracy comparison across CNN embedding types}
\label{tab:news_cnn_comparison}
\end{table}

BERT embeddings achieved the best results, as they provide rich contextual representations learned from large-scale pre-training. Random embeddings performed better than Word2Vec, likely because they can be fine-tuned specifically for the classification task during training, while Word2Vec embeddings are fixed and may not capture task-specific patterns as effectively.

\section{Summary}

\end{document}

